\documentclass[10pt,a4paper]{amsart}
\usepackage[margin=19mm]{geometry}
\usepackage{amsmath,amssymb,mathtools}
\usepackage[scr=rsfs,cal=euler]{mathalfa}
\usepackage{xfrac}
\usepackage[unicode,hidelinks,bookmarks=false]{hyperref}
\newcommand{\ie}{i.e.\ }
\newcommand{\cf}{cf.\ }
\newcommand{\resp}{respectively\ }
\newcommand{\Subsection}{\subsection}
\providecommand{\nobreakdash}{\nobreak\hbox{-}}
\setlength{\emergencystretch}{2em}
\multlinegap0pt
\usepackage{bm}
\usepackage{bbm}
\usepackage{dsfont}
\usepackage{xspace}
\usepackage{cancel}
\usepackage{mathtools}
\usepackage{amsthm}
\makeatletter
\@ifundefined{lemma}{\newtheorem{lemma}{Lemma}}{}
\makeatother
\title[On the pair correlation statistics for determinantal point p]{On the pair correlation statistics for determinantal point processes on the sphere}
\author{Maryna Manskova}
\date{Source: arXiv:2508.17952v2 (CC BY 4.0)}
\begin{document}
\maketitle

\noindent\emph{Source and licence.} On the pair correlation statistics for determinantal point processes on the sphere. Version arXiv:2508.17952v2 (CC BY 4.0), distributed under the Creative Commons Attribution 4.0 International licence, as recorded in the arXiv licence metadata for this exact version and shown on the abstract page https://arxiv.org/abs/2508.17952v2. Full rights notice: licenses/arxiv-2508.17952-CC-BY-4.0.txt.\par\smallskip

\noindent\textit{Editorial scope.} Counted unit: the Lemma whose source label is \texttt{total perimeter} in the article (source0.tex lines 580--583, source label \texttt{total perimeter}), delivered together with the article's own complete original proof of it (source0.tex lines 584--644, one \texttt{proof} environment of 61 lines). No mathematics is written, repaired or re-derived: statement and proof are the article's own text line for line. Required context retained uncounted: def y\_i (lines 557--559); def theta (lines 563--565). Every definition, display and label of those lines is retained unchanged. Omitted from this package and not redistributed: 1--556, 560--562, 566--579, 645--814. Nothing inside a retained excerpt is omitted or altered: every retained body is delivered exactly as the article's lines are written, so no omission receipt of any kind accompanies this package. Citations: the retained excerpts carry no citation command at all, so no bibliography entry of the article is transcribed and no reference is redistributed. Reference adaptation: none was needed and none was made; every reference inside the delivered unit resolves inside it, so no unresolved reference or citation is produced. Printed slips kept literal: no printed word, symbol, hypothesis or number of the counted statement or of its proof was altered. Layout and class: the article's own class is \texttt{amsart} with packages including amsfonts, amsthm, amssymb, amsmath, comment, a4wide, dsfont, bbm, xparse, mathrsfs, mathtools, enumitem, color, graphicx; the wrapper is the lane's pinned \texttt{amsart} editorial wrapper at 10pt on A4, which restates the theorem-like environments the retained text needs and adds the article's own macro definitions that the retained text needs (none), with extra wrapper packages (bm, bbm, dsfont, xspace, cancel, mathtools). The delivered numbering is the unit's own. Assets: no retained span contains a figure environment, a graphics-inclusion command, a PDF-page inclusion, a table or a TikZ picture; no image file is copied into this package, the package declares no asset dependency and no image file is redistributed with it. Licence: version arXiv:2508.17952v2 of this article is distributed under the Creative Commons Attribution 4.0 International licence, as recorded in the arXiv licence metadata for this exact version and shown on the abstract page https://arxiv.org/abs/2508.17952v2. A full rights notice is delivered at licenses/arxiv-2508.17952-CC-BY-4.0.txt. Characters: every retained excerpt is pure ASCII, so the delivered engine, XeLaTeX with its default Latin Modern text font, reports no missing character.
\begin{align}\label{def y_i}
    y_i=\frac{\sigma(C(\theta_{F,i+1}))-\sigma(C(\theta_{F,i}))}{1/N}=\frac{\cos\theta_{F,i}-\cos\theta_{F,i+1}}{2/N}.
\end{align}

\begin{align}\label{def theta}
    \sigma(C(\theta_i))=\frac{1}{N}\sum_{j=0}^{i-1}m_j.
\end{align}

\begin{lemma}\label{total perimeter}
    The sum of the perimeters of all regions in EQ$(2,N)$ is given by 
    $$P(N)=8\sqrt{\pi}\cdot\sqrt{N}+O(1) \quad \text{as}\quad N\rightarrow\infty.$$
\end{lemma}

\begin{proof}
    The union of all the boundaries consists of $n+1$ circles and $N-2$ arcs. The length of the circle at the colatitude $\theta_i$ is equal to $2\pi\sin\theta_i$. The length of one arc on the same level is $\theta_{i+1}-\theta_i$.
    The total length is 
\begin{align}\label{perimeter estimate}
    2\pi\sum\limits_{i=1}^{n+1}\sin\theta_i+\sum\limits_{i=1}^nm_i(\theta_{i+1}-\theta_i).
\end{align}
From the algorithm, it is known that $\theta_1=\theta_{F,1}=\theta_c$ and  $\theta_{n+1}=\theta_{F,n+1}=\pi-\theta_c$. Note that by (\ref{def y_i})  and (\ref{def theta}) we have 
$$\cos\theta_{i}=1-\frac{2}{N}\sum\limits_{j=0}^{i-1}m_j\quad \text{and} \quad \cos\theta_{F,i}=1-\frac{2}{N}\sum\limits_{j=0}^{i-1}y_j.$$
By the rounding procedure for $m_j$'s and $y_j$'s, 
$$\sum\limits_{j=0}^{i-2}y_j<\sum\limits_{j=0}^{i-1}y_j-\frac{1}{2}\leq \sum\limits_{j=0}^{i-1}m_j\leq \sum\limits_{j=0}^{i-1}y_j+\frac{1}{2}<\sum\limits_{j=0}^{i}y_j$$
which gives us 
$$\cos\theta_{F,i+1}<\cos\theta_{F,i}-\frac{1}{N}\leq \cos\theta_i\leq \cos \theta_{F,i}+\frac{1}{N}<\cos \theta_{F,i-1}$$
for $i=2,\dots ,n.$ It can be rewritten as 
$$|\cos\theta_i-\cos\theta_{F,i}|=\left|2\sin\tfrac{\theta_{F,i}+\theta_i}{2}\sin\tfrac{\theta_{F,i}-\theta_i}{2}\right|\leq \frac{1}{N},$$
Note that $\tfrac{\theta_{F,i}+\theta_i}{2}\in(\theta_{F,i-1},\theta_{F,i+1})\subseteq(\theta_c,\pi -\theta_c)$. As $\sin \tfrac{\theta_{F,i}+\theta_i}{2}\neq 0$, this gives us 
$$|\sin\theta_i-\sin\theta_{F,i}|=\left|2\cos\tfrac{\theta_{F,i}+\theta_i}{2}\sin\tfrac{\theta_{F,i}-\theta_i}{2}\right|\leq \frac{1}{N\sin\tfrac{\theta_{F,i}+\theta_i}{2}}\leq\frac{1}{\sqrt N}.$$
We can estimate the first sum in (\ref{perimeter estimate}),
\begin{align*}
    \left|\sum\limits_{i=1}^{n+1}\sin\theta_i-\sum\limits_{i=1}^{n+1}\sin\theta_{F,i}\right| &\leq \sum\limits_{i=2}^n\left|\sin\theta_i-\sin\theta_{F,i}\right|\leq \frac{n-1}{\sqrt N}=O(1) \quad \text{as}\quad N\rightarrow\infty.
\end{align*}
Using summation formulas for $\sin(kx)$ and $\cos(kx)$, we obtain 
\begin{align*}
    \sum\limits_{i=1}^{n+1}\sin\theta_{F,i}&=\sum\limits_{i=1}^{n+1}\sin(\theta_c+(i-1)\delta_F)\\ &=\sin\theta_c+\sin\theta_c\sum\limits_{k=1}^n \cos (k\delta_F)+\cos\theta_c\sum\limits_{k=1}^n\sin(k\delta_F)
    \\ &= \sin\theta_c+\sin\theta_c\cdot \frac{\sin\frac{n\delta_F}{2}\cos\frac{(n+1)\delta_F}{2}}{\sin\frac{\delta_F}{2}}+\cos\theta_c\cdot \frac{\sin\frac{n\delta_F}{2}\sin\frac{(n+1)\delta_F}{2}}{\sin\frac{\delta_F}{2}}.
\end{align*}
Using the fact that $n\delta_F=\pi-2\theta_c$, we get a simplified expression for the sum,
\begin{align*}
      \sum\limits_{i=1}^{n+1}\sin\theta_{F,i}&=\sin\theta_c+\sin\theta_c\cdot \frac{\cos\theta_c\sin(\theta_c-\frac{\delta_F}{2})}{\sin\frac{\delta_F}{2}}+\cos\theta_c\cdot \frac{\cos\theta_c\cos(\theta_c-\frac{\delta_F}{2})}{\sin\frac{\delta_F}{2}}
      \\ &=\sin\theta_c+\cos\theta_c\cot\frac{\delta_F}{2}
      \\ &=\sqrt{\frac{N}{\pi}}+O(1/\sqrt{N})\quad \text{as}\quad N\rightarrow\infty.
\end{align*}
For the second sum in (\ref{perimeter estimate}) we have
\begin{align*}
    \sum\limits_{i=1}^nm_i(\theta_{i+1}-\theta_i)&=\sum_{i=1}^nm_i-(\theta_{F,i+1}-\theta_{F,i}+\alpha_{i+1}-\alpha_i)
    \\&=\sum\limits_{i=1}^nm_i\delta_F+\sum\limits_{i=1}^nm_i(\alpha_{i+1}-\alpha_i)
    \\&=S_{1}+S_2,
\end{align*}
where $\alpha_i=\theta_i-\theta_{F,i}$. The main term is
$$S_1=\delta_F\sum\limits_{i=1}^nm_i=\left(\delta_I+O(1/N)\right)(N-2)=2\sqrt{\pi}\cdot\sqrt{N}+O(1).$$
Recall that $\alpha_1=\alpha_{n+1}=0$ and
\begin{equation*}
    m_{i}=\frac{\cos\theta_i-\cos\theta_{i+1}}{2/N}=N\sin\tfrac{\theta_i+\theta_{i+1}}{2}\sin\tfrac{\theta_{i+1}-\theta_{i}}{2}\leq2\sqrt N\sin \tfrac{\theta_i+\theta_{i+1}}{2}.
\end{equation*}
It implies that 
\begin{equation*}
    |m_{i}-m_{i-1}|\leq 4\sqrt{N}\sin\left(\tfrac{\theta_i}{2}+\tfrac{\theta_{i-1}+\theta_{i+1}}{4}\right).
\end{equation*}
The second sum is
\begin{align*}
    |S_2|=\left|\sum_{i=2}^n(m_i-m_{i-1})\alpha_i\right|&\leq\frac{1}{\sqrt{N}} \sum\limits_{i=2}^n\frac{\sin\left(\tfrac{\theta_i}{2}+\tfrac{\theta_{i-1}+\theta_{i+1}}{4}\right)}{\sin\tfrac{\theta_{i}+\theta_{F,i}}{2}}
    \\&\leq \frac{1}{\sqrt{N}} \sum\limits_{i=2}^n\left(1+\tfrac{1}{\sqrt{N}}\left|\cot\tfrac{\theta_{i}+\theta_{F,i}}{2}\right|\right)
    \\ &=\frac{1}{N}\sum\limits_{i=2}^n\left|\cot\tfrac{\theta_{i}+\theta_{F,i}}{2}\right|+O(1)\quad \text{as}\quad N\rightarrow\infty.
\end{align*}
Applying Koksma's inequality for $n-1$ points $x_i=\frac{\theta_{F,i}+\theta_i}{2}$ on $(\theta_c,\pi-\theta_c)$ and the function $f(x)=|\cot x|$, we estimate the error term by
\begin{align*}
   |S_2| &\leq \frac{n-1}{N}\int\limits_{\theta_c}^{\pi-\theta_c}|\cot(x)|\mathrm{d}x+\frac{n-1}{N}\cdot 2\cot \theta_c\cdot \frac{2}{n-1}+O(1)
    \\ &\leq\frac{1}{\sqrt{N}}\cdot 2\log\left(\frac{1}{\sin\theta_c}\right)+\frac{4\cot \theta_c}{N}+O(1)
    \\ & =O\left(1\right)\quad \text{as}\quad N\rightarrow\infty.
\end{align*}
Combining all together we get the desired result.
\end{proof} 

\end{document}
